\subsection{Subalgebras of an etale algebra}

PROPOSITION 3. --- Let A be an etale algebra over K. There exist only a finite number of subalgebras and ideals of A. Moreover, every extension of K which diagonalizes A also diagonalizes every subalgebra and every quotient algebra of A; in particular these algebras are etale.

It suffices to show that an algebra \(K^{n}\) has only a finite number of subalgebras and ideals, and that the subalgebras and quotient algebras of \(K^{n}\) are diagonalizable. We denote by \((\varepsilon_{1},\ldots,\varepsilon_{n})\) the canonical basis of \(K^{n}\) .

Let A be a subalgebra of \(K^{n}\) and let \(v_{1}, \ldots, v_{n}\) be the restrictions to A of the n projections \(K^{n} \rightarrow K\) . Since the intersection of the kernels of the \(v_{i}\) is clearly 0, the \(v_{i}\) generate the vector K-space dual to A (II, p.~302, Cor. 1); hence the K-algebra A is diagonalizable (V, p.~29, Prop. 1).

For every subset I of \(\{1, 2, \ldots, n\}\) put \(\varepsilon_{I} = \sum_{i \in I} \varepsilon_{i}\) . It is clear that the elements

\(\varepsilon_{I}\) are idempotents of \(K^{n}\) ; we have \(\varepsilon_{I}=0\) if and only if I is empty, and \(\varepsilon_{I}\varepsilon_{J}=\varepsilon_{I\cap J}\) . By what has been said, every subalgebra A of \(K^{n}\) is diagonalizable; by condition b) of Prop. 1 every subalgebra A of \(K^{n}\) therefore admits a basis \((\varepsilon_{I_{1}},\ldots,\varepsilon_{I_{p}})\) , where \((I_{1},\ldots,I_{p})\) is a partition of \(\{1,2,\ldots,n\}\) , and there are only a finite number of such subalgebras.

For every subset I of \(\{1,2,\ldots,n\}\) let \(a_{I}\) be the vector subspace of \(K^{n}\) having as basis the idempotents \(\varepsilon_{i}\) for \(i\in I\) ; then it is clear that \(a_{I}\) is an ideal of \(K^{n}\) ; moreover if \(J=\{1,2,\ldots,n\}-I\) , then the residue classes \(\overline{\varepsilon}_{j}\) of \(\varepsilon_{j} \mod a_{I}\) for \(j\in J\) form a basis of \(K^{n}/a_{I}\) . We have \(\overline{\varepsilon}_{j}^{2}=\overline{\varepsilon}_{j}\) and \(\overline{\varepsilon}_{j}\overline{\varepsilon}_{k}=0\) if \(j\neq k\) , hence the algebra \(K^{n}/a_{I}\) is diagonalizable, by Prop. 1 of V, p.~29.

It remains to show that every ideal of \(K^{n}\) is of the form \(a_{I}\) . Let I be the set of integers i such that \(1 \leqslant i \leqslant n\) and \(\varepsilon_{i} \in a\) , then \(a_{I} \subset a\) . Let \(x = x_{1}\varepsilon_{1} + \cdots + x_{n}\varepsilon_{n}\) be an element of a (with \(x_{1}, \ldots, x_{n}\) in K) and let i be in \(\{1, 2, \ldots, n\} - I\) . We have \(x_{i}\varepsilon_{i} = x\varepsilon_{i} \in a\) , and \(\varepsilon_{i} \notin a\) , whence \(x_{i} = 0\) . Thus \(x = \sum_{i \in I} x_{i}\varepsilon_{i}\) and this shows that

\(x \in a_{1}\) . We have now shown \(a \subset a_{1}\) , whence finally \(a = a_{1}\) .

COROLLARY. --- Let \(A_{1}, \ldots, A_{m}\) be algebras over K and \(A = A_{1} \times \cdots \times A_{m}\) . For A to be etale it is necessary and sufficient that \(A_{1}, \ldots, A_{m}\) are etale.

Suppose that A is etale; each of the algebras \(A_{1}, \ldots, A_{m}\) is isomorphic to a quotient of A, and hence is etale by Prop. 3. Conversely, every extension of K which diagonalizes \(A_{1}, \ldots, A_{m}\) clearly diagonalizes A.

